\documentclass[12pt,oneside]{report}

\usepackage[T1]{fontenc}
\usepackage[utf8]{inputenc}
\usepackage[brazil]{babel}
\usepackage{lmodern}
\usepackage{microtype}
\usepackage{indentfirst}
\usepackage[a4paper,margin=2.5cm]{geometry}

\usepackage{array}
\usepackage{longtable}
\usepackage{booktabs}

\usepackage{hyperref}
\usepackage{xurl}
\hypersetup{hidelinks}

\setcounter{secnumdepth}{3}
\setcounter{tocdepth}{3}

\begin{document}

\chapter{Plataforma de Observabilidade e Experimenta\c{c}\~ao em SDN}\label{ch:plataforma-observabilidade-sdn}

\section{Motiva\c{c}\~ao e objetivo}
Redes definidas por software (\textit{Software-Defined Networking} -- SDN) deslocam a complexidade operacional para o plano de controle e para a programa\c{c}\~ao do \textit{dataplane}. Essa mudan\c{c}a traz flexibilidade, mas amplifica a necessidade de observa\c{c}\~ao sistem\'atica: falhas e degrada\c{c}\~oes podem emergir em diferentes camadas (do enlace ao servi\c{c}o fim-a-fim) e, em SDN, tamb\'em no acoplamento entre controlador e switches.

Este cap\'itulo descreve a arquitetura e o funcionamento da plataforma \texttt{profissa-sdn-framework} (diret\'orio \texttt{platform/}), concebida para: (i) coletar, normalizar e persistir m\'etricas de m\'ultiplas camadas; (ii) oferecer uma interface para inspe\c{c}\~ao em tempo real via \textit{Server-Sent Events} (SSE) e consultas REST; e (iii) estruturar experimentos como registros cient\'ificos reprodut\'iveis (topologia, tr\'afego, plano de m\'etricas, execu\c{c}\~oes e artefatos).

\section{Vis\~ao geral da arquitetura}
A plataforma adota uma arquitetura cliente--servidor com persist\^encia em arquivos:

\begin{itemize}
\item \textbf{Backend} (FastAPI + Pydantic): implementa API REST e SSE, orquestra a coleta (real ou sint\'etica), aplica configura\c{c}\~oes ativas, realiza persist\^encia e exp\~oe cat\'alogo de m\'etricas. O n\'ucleo est\'a em \texttt{platform/backend/api/app.py}.
\item \textbf{Frontend} (Vite + React/TypeScript): consome SSE para telemetria ao vivo e REST para opera\c{c}\~oes de invent\'ario/configura\c{c}\~ao/experimentos. Rotas de UI est\~ao em \texttt{platform/frontend/src/pages/}.
\item \textbf{Persist\^encia por arquivos}: artefatos operacionais em \texttt{temp/} (configs, runs, m\'etricas por run) e telemetria em \texttt{raw/} (arquivos \texttt{metrics\_\{layer\}.jsonl}).
\end{itemize}

\subsection{Fluxo de dados (telemetria ao vivo e registro)}
A telemetria segue um padr\~ao \textit{append-only}:

\begin{enumerate}
\item Um processo de coleta gera amostras normalizadas (\texttt{MetricSample}).
\item As amostras s\~ao armazenadas no \textit{buffer} em mem\'oria e persistidas em \texttt{raw/} (JSONL por camada).
\item Eventos SSE (endpoint \texttt{/stream/events}) transportam um \textit{payload} JSON para o frontend, contendo \texttt{generated\_at}, \texttt{topology} e um lote de \texttt{metrics}.
\item Para fins de reprodutibilidade, execu\c{c}\~oes de experimento (\texttt{ExperimentRun}) registram um \textit{snapshot}: coleta \`unica de amostras e persist\^encia do conjunto como \texttt{MetricRecord} associado ao \texttt{run\_id}.
\end{enumerate}

No frontend, o consumo SSE est\'a centralizado no \textit{hook} \texttt{useEventStream} (\texttt{platform/frontend/src/hooks/useEventStream.ts}), que converte mensagens em \texttt{EventPayload} e atualiza o estado de telas que operam em tempo real (Dashboard, Monitor, Topologia, Lab).

\section{Modelo de dados e normaliza\c{c}\~ao}
A tipagem can\^onica \`e definida em \texttt{platform/backend/schemas.py}. Os principais conceitos s\~ao:

\begin{itemize}
\item \textbf{Camadas}: \texttt{MetricLayer} organiza m\'etricas em \texttt{service}, \texttt{physical}, \texttt{link}, \texttt{network}, \texttt{transport}, \texttt{session}, \texttt{presentation}, \texttt{application}, \texttt{control} e \texttt{dataplane}.
\item \textbf{Defini\c{c}\~oes}: \texttt{MetricDefinition} descreve nome, descri\c{c}\~ao, unidade, \textit{kind} (ex.: \texttt{gauge}, \texttt{counter}, \texttt{summary}, \texttt{derived}) e dimens\~oes (r\'otulos) esperadas.
\item \textbf{Amostras}: \texttt{MetricSample} representa um ponto de telemetria com \texttt{timestamp}, \texttt{layer}, \texttt{node}, \texttt{metric} e \texttt{value}, opcionalmente com \texttt{labels} e \texttt{details}.
\item \textbf{Registros}: \texttt{MetricRecord} agrega a amostra com metadados de execu\c{c}\~ao (ex.: \texttt{run\_id}, \texttt{experiment\_id}, \texttt{topology\_id}) quando a coleta \`e feita no contexto de experimentos.
\end{itemize}

A normaliza\c{c}\~ao (incluindo garantia de r\'otulos m\'inimos e coer\^encia de formato) \`e implementada em \texttt{platform/backend/core/metrics.py}.

\section{Cat\'alogo de m\'etricas (completude e rastreabilidade)}
O cat\'alogo can\^onico da plataforma est\'a em \texttt{platform/backend/core/metric\_catalog.py} e cont\'em 267 defini\c{c}\~oes. Cada defini\c{c}\~ao especifica: camada, unidade, \textit{kind}, categoria e dimens\~oes (r\'otulos). Esse cat\'alogo funciona como contrato entre coleta, persist\^encia e visualiza\c{c}\~ao: o backend exp\~oe \texttt{/metrics/definitions} e o frontend utiliza tais defini\c{c}\~oes para filtros, agrupamentos e contextualiza\c{c}\~ao de valores.

A Tabela~\ref{tab:metric-catalog} inclui o cat\'alogo completo (gerado diretamente a partir do reposit\'orio para evitar lacunas). Para compila\c{c}\~ao, recomenda-se suporte a \texttt{longtable} e \texttt{booktabs}.

\begingroup
\footnotesize
\setlength{\tabcolsep}{4pt}
\renewcommand{\arraystretch}{1.08}
\begin{longtable}{p{0.22\textwidth}p{0.10\textwidth}p{0.07\textwidth}p{0.09\textwidth}p{0.10\textwidth}p{0.16\textwidth}p{0.22\textwidth}}
\caption{Cat\'alogo de m\'etricas da plataforma (267 defini\c{c}\~oes).}\label{tab:metric-catalog}\\
\toprule
\textbf{Nome} & \textbf{Camada} & \textbf{Unid.} & \textbf{Kind} & \textbf{Categoria} & \textbf{Dimens\~oes} & \textbf{Descri\c{c}\~ao}\\
\midrule
\endfirsthead
\toprule
\textbf{Nome} & \textbf{Camada} & \textbf{Unid.} & \textbf{Kind} & \textbf{Categoria} & \textbf{Dimens\~oes} & \textbf{Descri\c{c}\~ao}\\
\midrule
\endhead
\midrule
\multicolumn{7}{r}{\footnotesize\emph{Continua na pr\'oxima p\'agina.}}\\
\endfoot
\bottomrule
\endlastfoot
% Auto-gerado a partir de platform.backend.core.metric_catalog (total=267)
\texttt{e2e\_path\_availability\_pct} & service & \% & ratio & e2e & src, dst, service & Disponibilidade do caminho (sucesso de probes sintéticos) \\
\texttt{e2e\_latency\_ms\_p50} & service & ms & summary & e2e & src, dst, service & Latência E2E p50 \\
\texttt{e2e\_latency\_ms\_p95} & service & ms & summary & e2e & src, dst, service & Latência E2E p95 \\
\texttt{e2e\_latency\_ms\_p99} & service & ms & summary & e2e & src, dst, service & Latência E2E p99 \\
\texttt{e2e\_jitter\_ms\_p95} & service & ms & summary & e2e & src, dst, service & Jitter E2E p95 \\
\texttt{e2e\_loss\_pct} & service & \% & ratio & e2e & src, dst, service & Perda E2E \\
\texttt{path\_changes\_total} & service & count & counter & e2e & src, dst, service & Mudanças de rota (path changes / ECMP hash changes) \\
\texttt{mttd\_sec} & service & s & derived & reliability & service, incident\_id & Tempo médio de detecção (MTTD) \\
\texttt{mttr\_sec} & service & s & derived & reliability & service, incident\_id & Tempo médio de restauração (MTTR) \\
\texttt{worst\_path\_score} & service &  & derived & e2e & src, dst, service & Score derivado para rankear piores caminhos (loss/latency/jitter) \\
\texttt{if\_link\_up} & physical & bool & gauge & health & node, interface & Status do link (up=1/down=0) \\
\texttt{if\_flaps\_total} & physical & count & counter & health & node, interface & Contagem de flaps (up/down) \\
\texttt{if\_speed\_mbps} & physical & Mbps & gauge & capacity & node, interface & Velocidade negociada \\
\texttt{if\_duplex\_mismatch} & physical & bool & gauge & health & node, interface & Mismatch de duplex (1=sim) \\
\texttt{if\_in\_bps} & physical & bps & gauge & capacity & node, interface & Tráfego de entrada (bits/s) \\
\texttt{if\_out\_bps} & physical & bps & gauge & capacity & node, interface & Tráfego de saída (bits/s) \\
\texttt{if\_in\_bps\_p95} & physical & bps & derived & capacity & node, interface & Utilização de entrada p95 (bits/s) \\
\texttt{if\_in\_bps\_p99} & physical & bps & derived & capacity & node, interface & Utilização de entrada p99 (bits/s) \\
\texttt{if\_out\_bps\_p95} & physical & bps & derived & capacity & node, interface & Utilização de saída p95 (bits/s) \\
\texttt{if\_out\_bps\_p99} & physical & bps & derived & capacity & node, interface & Utilização de saída p99 (bits/s) \\
\texttt{if\_errors} & physical & count & counter & errors & node, interface & Interface errors \\
\texttt{if\_discards} & physical & count & counter & errors & node, interface & Interface discards \\
\texttt{if\_crc\_errors} & physical & count & counter & errors & node, interface & CRC errors \\
\texttt{if\_fcs\_errors\_total} & physical & count & counter & errors & node, interface & FCS errors \\
\texttt{if\_alignment\_errors\_total} & physical & count & counter & errors & node, interface & Alignment errors \\
\texttt{if\_symbol\_errors\_total} & physical & count & counter & errors & node, interface & Symbol errors \\
\texttt{if\_physical\_drops\_total} & physical & count & counter & errors & node, interface & Drops físicos (quando possível separar de congestionamento) \\
\texttt{if\_optic\_rx\_dbm} & physical & dBm & gauge & health & node, interface & Optical RX power \\
\texttt{if\_optic\_tx\_dbm} & physical & dBm & gauge & health & node, interface & Optical TX power \\
\texttt{if\_optic\_osnr\_db} & physical & dB & gauge & health & node, interface & Optical OSNR \\
\texttt{if\_optic\_bias\_ma} & physical & mA & gauge & health & node, interface & Optical bias current \\
\texttt{if\_optic\_temp\_c} & physical & C & gauge & health & node, interface & Optical module temperature \\
\texttt{if\_optic\_threshold\_violations\_total} & physical & count & counter & health & node, interface & Optical threshold violations \\
\texttt{if\_temp\_c} & physical & C & gauge & health & node, interface & Interface temperature \\
\texttt{if\_mtu\_effective\_bytes} & physical & bytes & gauge & health & node, interface & MTU efetivo (detecção de inconsistências) \\
\texttt{poe\_power\_w} & physical & W & gauge & capacity & node, interface & Consumo PoE \\
\texttt{poe\_events\_total} & physical & count & counter & health & node, interface & Eventos PoE \\
\texttt{poe\_budget\_w} & physical & W & gauge & capacity & node & Budget PoE \\
\texttt{mac\_table\_size} & link & count & gauge & l2 & node, vlan & Tamanho da MAC table \\
\texttt{mac\_churn\_rate} & link & count/s & derived & l2 & node, vlan & MAC churn rate \\
\texttt{mac\_moves\_total} & link & count & counter & l2 & node, vlan & MAC moves \\
\texttt{stp\_root\_changes\_total} & link & count & counter & l2 & node, instance & STP root changes \\
\texttt{stp\_topology\_changes\_total} & link & count & counter & l2 & node, instance & STP topology changes \\
\texttt{stp\_blocked\_ports} & link & count & gauge & l2 & node & STP blocked ports \\
\texttt{lacp\_members\_up} & link & count & gauge & l2 & node, portchannel & LACP members up \\
\texttt{lacp\_members\_down} & link & count & gauge & l2 & node, portchannel & LACP members down \\
\texttt{lacp\_imbalance\_pct} & link & \% & ratio & capacity & node, portchannel & LACP imbalance \\
\texttt{lacp\_renegotiations\_total} & link & count & counter & l2 & node, portchannel & LACP renegotiations \\
\texttt{vlan\_active\_count} & link & count & gauge & l2 & node, trunk & VLANs ativas por trunk \\
\texttt{vlan\_allowed\_mismatch\_total} & link & count & counter & l2 & node, trunk & Inconsistência de VLANs allowed \\
\texttt{vlan\_native\_mismatch\_total} & link & count & counter & l2 & node, trunk & Inconsistência de VLAN native \\
\texttt{bum\_rate\_pps} & link & pps & gauge & capacity & node, vlan, type & Taxa de BUM (broadcast/multicast/unknown) \\
\texttt{bum\_peak\_pps} & link & pps & derived & capacity & node, vlan, type & Pico de BUM \\
\texttt{arp\_cache\_util\_pct} & link & \% & ratio & l2 & node, vrf & ARP cache utilization \\
\texttt{arp\_miss\_rate} & link & count/s & derived & l2 & node, vrf & ARP miss rate \\
\texttt{arp\_duplicates\_total} & link & count & counter & l2 & node, vrf & ARP duplicados \\
\texttt{arp\_gratuitous\_spikes\_total} & link & count & counter & l2 & node, vrf & Picos de gratuitous ARP \\
\texttt{nd\_cache\_util\_pct} & link & \% & ratio & l2 & node, vrf & ND (IPv6) cache utilization \\
\texttt{nd\_miss\_rate} & link & count/s & derived & l2 & node, vrf & ND miss rate \\
\texttt{nd\_duplicates\_total} & link & count & counter & l2 & node, vrf & ND duplicados \\
\texttt{lldp\_neighbors\_expected\_pct} & link & \% & ratio & l2 & node & LLDP/CDP expected vs real \\
\texttt{lldp\_drift\_total} & link & count & counter & l2 & node & LLDP/CDP drift events \\
\texttt{microburst\_queue\_occupancy\_pct\_peak} & link & \% & derived & capacity & node, interface, queue & Microbursts: pico de fila/ocupação \\
\texttt{link\_util\_pct} & link & \% & gauge & capacity & node, interface & Link utilization \\
\texttt{link\_loss\_pct} & link & \% & ratio & errors & node, interface & Link loss \\
\texttt{link\_jitter\_ms} & link & ms & gauge & performance & node, interface & Link jitter \\
\texttt{queue\_occupancy\_pct} & link & \% & gauge & capacity & node, interface, queue, qos\_class & Queue occupancy \\
\texttt{queue\_drops} & link & count & counter & errors & node, interface, queue, qos\_class & Queue drops \\
\texttt{ecn\_marked\_pct} & link & \% & ratio & capacity & node, interface, qos\_class & ECN marked traffic \\
\texttt{prefix\_reachability\_pct} & network & \% & ratio & l3 & prefix, vrf, tenant, src, dst & Reachability por prefixo \\
\texttt{routing\_convergence\_ms} & network & ms & derived & l3 & protocol, vrf & Tempo de convergência de roteamento \\
\texttt{route\_count} & network & count & gauge & l3 & protocol, afi, safi, vrf & Route count (total/por AFI/SAFI/VRF) \\
\texttt{route\_churn\_rate} & network & count/s & derived & l3 & protocol, vrf & Route churn rate \\
\texttt{bgp\_session\_up} & network & bool & gauge & routing & peer, vrf, afi, safi & BGP session up (1/0) \\
\texttt{bgp\_session\_flaps\_total} & network & count & counter & routing & peer, vrf & BGP session flaps \\
\texttt{bgp\_updates\_total} & network & count & counter & routing & peer, vrf, afi, safi & BGP updates \\
\texttt{bgp\_withdraws\_total} & network & count & counter & routing & peer, vrf, afi, safi & BGP withdraws \\
\texttt{bgp\_prefix\_limit\_hits\_total} & network & count & counter & routing & peer, vrf & BGP prefix limit hits \\
\texttt{bgp\_as\_path\_changes\_total} & network & count & counter & routing & prefix, peer, vrf & BGP AS-path changes \\
\texttt{ospf\_session\_up} & network & bool & gauge & routing & neighbor, area, vrf & OSPF adjacency up (1/0) \\
\texttt{ospf\_flaps\_total} & network & count & counter & routing & neighbor, area, vrf & OSPF adjacency flaps \\
\texttt{isis\_session\_up} & network & bool & gauge & routing & neighbor, level, vrf & IS-IS adjacency up (1/0) \\
\texttt{isis\_flaps\_total} & network & count & counter & routing & neighbor, level, vrf & IS-IS adjacency flaps \\
\texttt{icmp\_unreachable\_total} & network & count & counter & errors & src, dst, code & ICMP unreachable \\
\texttt{icmp\_frag\_needed\_total} & network & count & counter & errors & src, dst & ICMP frag needed (PMTUD) \\
\texttt{ip\_fragments\_total} & network & count & counter & errors & src, dst & IP fragments \\
\texttt{ip\_fragment\_drops\_total} & network & count & counter & errors & node & Drops por MTU/fragmentação \\
\texttt{dscp\_packets\_total} & network & count & counter & qos & dscp, node, interface & Distribuição DSCP (contagem por classe) \\
\texttt{dscp\_remark\_events\_total} & network & count & counter & qos & from, to, policy & Eventos de remarking DSCP \\
\texttt{acl\_drops\_total} & network & count & counter & security & rule\_id, node & ACL/Policy drops por regra \\
\texttt{tunnel\_up} & network & bool & gauge & overlay & tunnel, type, vrf & Tunnel up/down \\
\texttt{tunnel\_keepalive\_loss\_total} & network & count & counter & overlay & tunnel, type, vrf & Tunnel keepalive loss \\
\texttt{tunnel\_encap\_errors\_total} & network & count & counter & overlay & tunnel, type & Encapsulation errors \\
\texttt{tunnel\_decap\_errors\_total} & network & count & counter & overlay & tunnel, type & Decapsulation errors \\
\texttt{vxlan\_vni\_mapping\_ok} & network & bool & gauge & overlay & vni, vrf & VNI/VRF mapping health \\
\texttt{latency\_ms} & network & ms & gauge & performance & src, dst & Latency \\
\texttt{packet\_loss\_pct} & network & \% & ratio & performance & src, dst & Packet loss \\
\texttt{jitter\_ms} & network & ms & gauge & performance & src, dst & Jitter \\
\texttt{ttl\_expired} & network & count & counter & errors & src, dst & TTL expired events \\
\texttt{routes\_count} & network & count & gauge & routing & node, vrf & Route entries \\
\texttt{arp\_entries} & network & count & gauge & l3 & node, vrf & ARP/ND cache entries \\
\texttt{tcp\_retrans\_pct} & transport & \% & ratio & transport & src, dst, flow & TCP retransmissions (\%) \\
\texttt{tcp\_retrans\_total} & transport & count & counter & transport & src, dst, flow & TCP retransmissions (count) \\
\texttt{tcp\_fast\_retrans\_total} & transport & count & counter & transport & src, dst, flow & TCP fast retransmit \\
\texttt{tcp\_dup\_acks\_total} & transport & count & counter & transport & src, dst, flow & TCP duplicate ACKs \\
\texttt{tcp\_handshake\_success\_pct} & transport & \% & ratio & transport & dst, port & TCP handshake success rate \\
\texttt{tcp\_handshake\_latency\_ms} & transport & ms & gauge & transport & dst, port & TCP handshake latency \\
\texttt{tcp\_resets\_total} & transport & count & counter & transport & src, dst, flow & Resets (RST) \\
\texttt{tcp\_out\_of\_order\_total} & transport & count & counter & transport & src, dst, flow & TCP out-of-order \\
\texttt{tcp\_zero\_window\_total} & transport & count & counter & transport & src, dst, flow & TCP zero window \\
\texttt{tcp\_window\_scaling\_anomalies\_total} & transport & count & counter & transport & src, dst, flow & TCP window scaling anomalies \\
\texttt{udp\_loss\_pct} & transport & \% & ratio & transport & src, dst, flow & UDP loss \\
\texttt{udp\_jitter\_ms} & transport & ms & gauge & transport & src, dst, flow & UDP jitter \\
\texttt{nat\_table\_util\_pct} & transport & \% & ratio & nat & node, nat\_pool & NAT table utilization \\
\texttt{nat\_port\_exhaustion\_events\_total} & transport & count & counter & nat & node, nat\_pool & NAT port exhaustion events \\
\texttt{nat\_translation\_failures\_total} & transport & count & counter & nat & node & NAT translation failures \\
\texttt{lb\_healthcheck\_success\_pct} & transport & \% & ratio & lb & lb, pool, backend & Load balancer health checks success \\
\texttt{lb\_connection\_rate} & transport & count/s & gauge & lb & lb, vip & Load balancer connection rate \\
\texttt{lb\_concurrent\_connections} & transport & count & gauge & lb & lb, vip & Load balancer concurrent connections \\
\texttt{lb\_backend\_connect\_errors\_total} & transport & count & counter & lb & lb, backend & LB backend connect errors \\
\texttt{lb\_backend\_timeouts\_total} & transport & count & counter & lb & lb, backend & LB backend timeouts \\
\texttt{throughput\_mbps} & transport & Mbps & gauge & performance & src, dst & Throughput \\
\texttt{tcp\_rtt\_ms} & transport & ms & gauge & transport & src, dst, flow & TCP RTT \\
\texttt{ss\_sockets\_total} & transport & count & gauge & host & node & Total sockets from ss -s \\
\texttt{ss\_tcp\_states} & transport & count & gauge & host & node & TCP states from ss -s \\
\texttt{ss\_udp\_sockets\_total} & transport & count & gauge & host & node & UDP sockets from ss -s \\
\texttt{ss\_tcp\_sockets\_total} & transport & count & gauge & host & node & TCP sockets from ss -s \\
\texttt{ss\_raw\_sockets\_total} & transport & count & gauge & host & node & RAW sockets from ss -s \\
\texttt{ss\_inet\_sockets\_total} & transport & count & gauge & host & node & INET sockets from ss -s \\
\texttt{ss\_frag\_sockets\_total} & transport & count & gauge & host & node & FRAG sockets from ss -s \\
\texttt{snmp\_ip\_in\_receives} & transport & count & counter & snmp & node & IP InReceives \\
\texttt{snmp\_ip\_in\_hdr\_errors} & transport & count & counter & snmp & node & IP InHdrErrors \\
\texttt{snmp\_ip\_in\_addr\_errors} & transport & count & counter & snmp & node & IP InAddrErrors \\
\texttt{snmp\_ip\_forw\_datagrams} & transport & count & counter & snmp & node & IP ForwDatagrams \\
\texttt{snmp\_ip\_in\_unknown\_protos} & transport & count & counter & snmp & node & IP InUnknownProtos \\
\texttt{snmp\_ip\_in\_discards} & transport & count & counter & snmp & node & IP InDiscards \\
\texttt{snmp\_ip\_in\_delivers} & transport & count & counter & snmp & node & IP InDelivers \\
\texttt{snmp\_ip\_out\_requests} & transport & count & counter & snmp & node & IP OutRequests \\
\texttt{snmp\_ip\_out\_discards} & transport & count & counter & snmp & node & IP OutDiscards \\
\texttt{snmp\_ip\_out\_no\_routes} & transport & count & counter & snmp & node & IP OutNoRoutes \\
\texttt{snmp\_ip\_reasm\_timeout} & transport & count & counter & snmp & node & IP ReasmTimeout \\
\texttt{snmp\_ip\_reasm\_reqds} & transport & count & counter & snmp & node & IP ReasmReqds \\
\texttt{snmp\_ip\_reasm\_oks} & transport & count & counter & snmp & node & IP ReasmOKs \\
\texttt{snmp\_ip\_reasm\_fails} & transport & count & counter & snmp & node & IP ReasmFails \\
\texttt{snmp\_ip\_frag\_oks} & transport & count & counter & snmp & node & IP FragOKs \\
\texttt{snmp\_ip\_frag\_fails} & transport & count & counter & snmp & node & IP FragFails \\
\texttt{snmp\_ip\_frag\_creates} & transport & count & counter & snmp & node & IP FragCreates \\
\texttt{snmp\_ip\_out\_transmits} & transport & count & counter & snmp & node & IP OutTransmits \\
\texttt{snmp\_icmp\_in\_msgs} & transport & count & counter & snmp & node & ICMP InMsgs \\
\texttt{snmp\_icmp\_out\_msgs} & transport & count & counter & snmp & node & ICMP OutMsgs \\
\texttt{snmp\_icmp\_in\_errors} & transport & count & counter & snmp & node & ICMP InErrors \\
\texttt{snmp\_icmp\_out\_errors} & transport & count & counter & snmp & node & ICMP OutErrors \\
\texttt{snmp\_icmp\_in\_dest\_unreachs} & transport & count & counter & snmp & node & ICMP InDestUnreachs \\
\texttt{snmp\_icmp\_out\_dest\_unreachs} & transport & count & counter & snmp & node & ICMP OutDestUnreachs \\
\texttt{snmp\_icmp\_in\_time\_excds} & transport & count & counter & snmp & node & ICMP InTimeExcds \\
\texttt{snmp\_icmp\_out\_time\_excds} & transport & count & counter & snmp & node & ICMP OutTimeExcds \\
\texttt{snmp\_icmp\_in\_echo\_reqs} & transport & count & counter & snmp & node & ICMP InEchos \\
\texttt{snmp\_icmp\_out\_echo\_reps} & transport & count & counter & snmp & node & ICMP OutEchoReps \\
\texttt{snmp\_tcp\_active\_opens} & transport & count & counter & snmp & node & TCP ActiveOpens \\
\texttt{snmp\_tcp\_passive\_opens} & transport & count & counter & snmp & node & TCP PassiveOpens \\
\texttt{snmp\_tcp\_attempt\_fails} & transport & count & counter & snmp & node & TCP AttemptFails \\
\texttt{snmp\_tcp\_estab\_resets} & transport & count & counter & snmp & node & TCP EstabResets \\
\texttt{snmp\_tcp\_curr\_estab} & transport & count & gauge & snmp & node & TCP CurrEstab \\
\texttt{snmp\_tcp\_in\_segs} & transport & count & counter & snmp & node & TCP InSegs \\
\texttt{snmp\_tcp\_out\_segs} & transport & count & counter & snmp & node & TCP OutSegs \\
\texttt{snmp\_tcp\_retrans\_segs} & transport & count & counter & snmp & node & TCP RetransSegs \\
\texttt{snmp\_tcp\_in\_errs} & transport & count & counter & snmp & node & TCP InErrs \\
\texttt{snmp\_tcp\_out\_rsts} & transport & count & counter & snmp & node & TCP OutRsts \\
\texttt{snmp\_tcp\_in\_csum\_errors} & transport & count & counter & snmp & node & TCP InCsumErrors \\
\texttt{snmp\_udp\_in\_datagrams} & transport & count & counter & snmp & node & UDP InDatagrams \\
\texttt{snmp\_udp\_no\_ports} & transport & count & counter & snmp & node & UDP NoPorts \\
\texttt{snmp\_udp\_in\_errors} & transport & count & counter & snmp & node & UDP InErrors \\
\texttt{snmp\_udp\_out\_datagrams} & transport & count & counter & snmp & node & UDP OutDatagrams \\
\texttt{snmp\_udp\_rcvbuf\_errors} & transport & count & counter & snmp & node & UDP RcvbufErrors \\
\texttt{snmp\_udp\_sndbuf\_errors} & transport & count & counter & snmp & node & UDP SndbufErrors \\
\texttt{snmp\_udp\_in\_csum\_errors} & transport & count & counter & snmp & node & UDP InCsumErrors \\
\texttt{snmp\_udp\_ignored\_multi} & transport & count & counter & snmp & node & UDP IgnoredMulti \\
\texttt{snmp\_udp\_mem\_errors} & transport & count & counter & snmp & node & UDP MemErrors \\
\texttt{tls\_session\_resumption\_pct} & session & \% & ratio & tls & sni, service & TLS session resumption rate \\
\texttt{tls\_renegotiations\_total} & session & count & counter & tls & sni, service & TLS renegotiations \\
\texttt{tls\_handshake\_failures\_total} & session & count & counter & tls & sni, cert & TLS handshake failures (por SNI/cert) \\
\texttt{vpn\_sessions\_active} & session & count & gauge & vpn & profile, tenant & VPN sessions active \\
\texttt{vpn\_session\_establish\_rate} & session & count/s & gauge & vpn & profile, tenant & VPN session establish rate \\
\texttt{vpn\_session\_failures\_total} & session & count & counter & vpn & profile, tenant, reason & VPN session failures \\
\texttt{aaa\_auth\_failures\_total} & session & count & counter & aaa & realm, reason & Falhas de autenticação AAA \\
\texttt{session\_idle\_timeouts\_total} & session & count & counter & session & policy, profile & Idle/timeout expirations \\
\texttt{fw\_state\_table\_util\_pct} & session & \% & ratio & security & node & Firewall state table usage \\
\texttt{fw\_state\_table\_drops\_total} & session & count & counter & security & node & Drops por overflow de state table \\
\texttt{tls\_handshake\_latency\_ms\_p95} & presentation & ms & summary & tls & sni, service & TLS handshake latency p95 \\
\texttt{tls\_handshake\_latency\_ms\_p99} & presentation & ms & summary & tls & sni, service & TLS handshake latency p99 \\
\texttt{tls\_version\_connections\_total} & presentation & count & counter & tls & version, sni & Distribuição de versões TLS \\
\texttt{tls\_cipher\_connections\_total} & presentation & count & counter & tls & cipher, sni & Distribuição de ciphers \\
\texttt{cert\_expiry\_days} & presentation & days & gauge & tls & cert, sni & Certificados: dias até expiração \\
\texttt{cert\_validation\_errors\_total} & presentation & count & counter & tls & cert, sni & Certificados: validation errors \\
\texttt{cert\_chain\_errors\_total} & presentation & count & counter & tls & cert, sni & Certificados: chain errors \\
\texttt{sni\_mismatch\_total} & presentation & count & counter & tls & sni, cert & SNI mismatch \\
\texttt{compression\_errors\_total} & presentation & count & counter & encoding & algo & Erro de compressão (gzip/br) \\
\texttt{encoding\_errors\_total} & presentation & count & counter & encoding & codec & Erro de encoding/codec \\
\texttt{http2\_negotiation\_success\_pct} & presentation & \% & ratio & http & service & HTTP/2 negotiation success \\
\texttt{http3\_quic\_negotiation\_success\_pct} & presentation & \% & ratio & http & service & HTTP/3 (QUIC) negotiation success \\
\texttt{quic\_fallback\_rate\_pct} & presentation & \% & ratio & http & service & HTTP/3 fallback rate \\
\texttt{app\_rps} & application & rps & gauge & app & service, route & Taxa de requisições (RPS/QPS) \\
\texttt{app\_latency\_ms\_p95} & application & ms & summary & app & service, route & Latência p95 (app) \\
\texttt{app\_latency\_ms\_p99} & application & ms & summary & app & service, route & Latência p99 (app) \\
\texttt{http\_responses\_total} & application & count & counter & http & service, route, code & Códigos HTTP (por classe/rota) \\
\texttt{http\_errors\_total} & application & count & counter & http & service, route, code & HTTP errors (4xx/5xx) \\
\texttt{dns\_qps} & application & qps & gauge & dns & server, zone & DNS QPS \\
\texttt{dns\_latency\_ms\_p95} & application & ms & summary & dns & server, zone & DNS latency p95 \\
\texttt{dns\_servfail\_total} & application & count & counter & dns & server, zone & DNS SERVFAIL \\
\texttt{dns\_nxdomain\_total} & application & count & counter & dns & server, zone & DNS NXDOMAIN \\
\texttt{dns\_cache\_hit\_ratio\_pct} & application & \% & ratio & dns & server & DNS cache hit ratio \\
\texttt{dhcp\_success\_pct} & application & \% & ratio & dhcp & scope, server & DHCP success rate (DORA) \\
\texttt{dhcp\_dora\_latency\_ms} & application & ms & gauge & dhcp & scope, server & DHCP DORA latency \\
\texttt{dhcp\_pool\_util\_pct} & application & \% & ratio & dhcp & scope & DHCP pool utilization \\
\texttt{dhcp\_conflicts\_total} & application & count & counter & dhcp & scope & DHCP conflicts \\
\texttt{ntp\_offset\_ms} & application & ms & gauge & ntp & server & NTP offset/drift \\
\texttt{ntp\_drift\_ppm} & application & ppm & gauge & ntp & server & NTP drift \\
\texttt{ntp\_stratum\_changes\_total} & application & count & counter & ntp & server & NTP stratum changes \\
\texttt{ntp\_reachability\_pct} & application & \% & ratio & ntp & server & NTP reachability \\
\texttt{http\_ttfb\_ms} & application & ms & gauge & http & service, route & HTTP time-to-first-byte \\
\texttt{http\_ttlb\_ms} & application & ms & gauge & http & service, route & HTTP time-to-last-byte \\
\texttt{http\_payload\_bytes\_p95} & application & bytes & summary & http & service, route & HTTP payload size p95 \\
\texttt{grpc\_status\_total} & application & count & counter & grpc & service, method, status & gRPC status codes \\
\texttt{grpc\_deadlines\_exceeded\_total} & application & count & counter & grpc & service, method & gRPC deadlines exceeded \\
\texttt{db\_connection\_errors\_total} & application & count & counter & db & db, service & Database connection errors (alto nível) \\
\texttt{db\_query\_latency\_ms\_p95} & application & ms & summary & db & db, service & Database query latency p95 (alto nível) \\
\texttt{controller\_latency\_ms} & control & ms & gauge & control & controller, endpoint & Controller latency \\
\texttt{controller\_conn\_ok} & control & bool & gauge & control & controller & Controller connectivity \\
\texttt{of\_channel\_reconnects} & control & count & counter & control & switch & OpenFlow reconnects \\
\texttt{flows\_installed} & control & count & counter & control & controller, switch & Flows installed \\
\texttt{flows\_removed} & control & count & counter & control & controller, switch & Flows removed \\
\texttt{cpu\_util\_pct} & control & \% & gauge & host & node & CPU utilization \\
\texttt{mem\_util\_pct} & control & \% & gauge & host & node & Memory utilization \\
\texttt{openflow\_flow\_packets} & dataplane & packets & counter & dataplane & switch, table, flow & Flow packets \\
\texttt{openflow\_flow\_bytes} & dataplane & bytes & counter & dataplane & switch, table, flow & Flow bytes \\
\texttt{table\_hits} & dataplane & count & counter & dataplane & switch, table & Table hits \\
\texttt{table\_misses} & dataplane & count & counter & dataplane & switch, table & Table misses \\
\texttt{port\_rx\_pkts} & dataplane & packets & counter & dataplane & switch, port & Port RX packets \\
\texttt{port\_tx\_pkts} & dataplane & packets & counter & dataplane & switch, port & Port TX packets \\
\texttt{port\_rx\_drops} & dataplane & packets & counter & dataplane & switch, port & Port RX drops \\
\texttt{port\_tx\_drops} & dataplane & packets & counter & dataplane & switch, port & Port TX drops \\
\texttt{flow\_bytes} & dataplane & bytes & counter & dataplane & src, dst, 5tuple & Flow byte counter \\
\texttt{flow\_packets} & dataplane & packets & counter & dataplane & src, dst, 5tuple & Flow packet counter \\
\texttt{flow\_duration\_sec} & dataplane & s & gauge & dataplane & src, dst, 5tuple & Flow duration \\
\texttt{interface\_util\_p95\_pct} & service & \% & derived & capacity & node, interface & Interface utilization p95 \\
\texttt{interface\_headroom\_pct} & service & \% & derived & capacity & node, interface & Interface headroom \\
\texttt{queue\_depth} & service & count & gauge & capacity & node, interface, queue, qos\_class & Queue depth \\
\texttt{buffer\_occupancy\_pct} & service & \% & gauge & capacity & node, asic & Buffer occupancy \\
\texttt{ecn\_marks\_total} & service & count & counter & capacity & node, interface, qos\_class & ECN marks \\
\texttt{elephant\_flows\_total} & service & count & derived & capacity & node, app & Elephant flows (classificação) \\
\texttt{mice\_flows\_total} & service & count & derived & capacity & node, app & Mice flows (classificação) \\
\texttt{hotspot\_score} & service &  & derived & capacity & node, interface, asic & Hotspot score (links/ASICs com maior utilização/drops) \\
\texttt{slo\_error\_budget\_remaining\_pct} & service & \% & derived & reliability & slo, service & Error budget restante do SLO \\
\texttt{incident\_rate} & service & count/s & derived & reliability & service & Taxa de incidentes \\
\texttt{flap\_storm\_events\_total} & service & count & counter & reliability & node, interface & Flap storms \\
\texttt{config\_drift\_score} & service &  & derived & reliability & node & Config drift vs golden config \\
\texttt{firewall\_denies\_total} & service & count & counter & security & policy, rule\_id, node & Firewall/IDS denies por policy \\
\texttt{scan\_spikes\_total} & service & count & counter & security & src, dst, port & Picos de scans \\
\texttt{ddos\_syn\_rate\_pps} & service & pps & gauge & security & dst, service & DDoS signal: SYN rate \\
\texttt{ddos\_udp\_flood\_pps} & service & pps & gauge & security & dst, service & DDoS signal: UDP flood indicators \\
\texttt{abnormal\_pps} & service & pps & derived & security & node, interface & PPS anormal \\
\texttt{anomaly\_new\_asns\_total} & service & count & counter & security & asn, prefix & Anomalias: novos ASNs \\
\texttt{anomaly\_new\_countries\_total} & service & count & counter & security & country, service & Anomalias: novos países \\
\texttt{anomaly\_new\_services\_total} & service & count & counter & security & service & Anomalias: novos serviços \\
\texttt{anomaly\_new\_ports\_total} & service & count & counter & security & port, proto & Anomalias: novas portas \\
\texttt{auth\_failures\_total} & service & count & counter & security & realm, profile & Falhas AAA/VPN (agregado) \\
\texttt{cert\_errors\_total} & service & count & counter & security & sni, cert & Cert errors (agregado) \\

\end{longtable}
\endgroup

\section{API: contratos, opera\c{c}\~oes e seguran\c{c}a}
O backend exp\~oe endpoints REST para invent\'ario (topologias, fluxos, configs), experimenta\c{c}\~ao (experimentos e execu\c{c}\~oes) e m\'etricas (cat\'alogo, ingest\~ao, consulta, exporta\c{c}\~ao), al\'em de SSE para telemetria ao vivo.

\subsection{Endpoints principais}
Os endpoints centrais (ver \texttt{platform/backend/api/app.py}) incluem:

\begin{itemize}
\item \textbf{Topologias}: \texttt{GET /topologies}, \texttt{GET /topologies/\{id\}}, \texttt{POST /topologies}, \texttt{PUT /topologies/\{id\}}, \texttt{DELETE /topologies/\{id\}}.
\item \textbf{Experimentos}: \texttt{GET /experiments}, \texttt{GET /experiments/\{id\}}, \texttt{POST /experiments}, \texttt{PUT /experiments/\{id\}}, \texttt{DELETE /experiments/\{id\}}.
\item \textbf{Execu\c{c}\~oes}: \texttt{POST /experiments/\{id\}/run?mode=synthetic|real}, \texttt{GET /experiments/\{id\}/runs}, \texttt{GET /runs/\{run\_id\}}, \texttt{GET /runs/\{run\_id\}/metrics}.
\item \textbf{M\'etricas}: \texttt{GET /metrics/definitions}, \texttt{POST /metrics/query}, \texttt{POST /metrics/samples}, \texttt{GET /metrics/latest}, \texttt{GET /metrics/export}, \texttt{POST /metrics/collect}.
\item \textbf{Fluxos}: \texttt{GET /flows}, \texttt{POST /flows}, \texttt{DELETE /flows/\{flow\_id\}}.
\item \textbf{Configura\c{c}\~oes}: \texttt{GET /configs}, \texttt{GET /configs/active}, \texttt{GET /configs/\{id\}}, \texttt{POST /configs}, \texttt{POST /configs/\{id\}/activate}, \texttt{DELETE /configs/\{id\}}.
\item \textbf{Streaming}: \texttt{GET /stream/events}.
\item \textbf{Sa\'ude}: \texttt{GET /health}.
\end{itemize}

\subsection{Autentica\c{c}\~ao (API key)}
Opera\c{c}\~oes mut\'aveis (ex.: criar topologia/experimento, iniciar execu\c{c}\~ao, salvar/ativar configs) s\~ao protegidas por uma chave via cabe\c{c}alho \texttt{X-API-Key}, exigida no backend (depend\^encia \texttt{require\_api\_key}). No frontend, as telas de experimentos e configura\c{c}\~oes usam \texttt{VITE\_API\_KEY} para montar cabe\c{c}alhos de requisi\c{c}\~ao.

\section{Coleta de m\'etricas: modos real e sint\'etico}
A plataforma implementa dois modos de coleta:

\begin{itemize}
\item \textbf{Real}: utiliza comandos e fontes do sistema/ambiente (ex.: \texttt{docker exec}, \texttt{ping}, \texttt{ip -s link}, \texttt{/proc/net/snmp}, \texttt{ovs-ofctl}) para extrair contadores, lat\^encia/perda e sinais do plano de controle/dataplane.
\item \textbf{Sint\'etico}: gera amostras de forma determin\'istica/aleat\'oria para cobrir o cat\'alogo mesmo quando o ambiente real n\~ao est\'a dispon\'ivel, permitindo valida\c{c}\~ao de pipeline, dashboards e filtros.
\end{itemize}

Ambos os modos produzem \texttt{MetricSample} com labels de contexto (em particular \texttt{experiment\_id}, \texttt{run\_id} e \texttt{topology\_id} quando aplic\'avel), permitindo correla\c{c}\~ao e rastreabilidade de resultados.

\section{Interface (frontend): descri\c{c}\~ao por p\'agina}
As p\'aginas do frontend materializam diferentes vistas do mesmo \textit{data plane} de telemetria. As principais rotas est\~ao em \texttt{platform/frontend/src/main.tsx} e as telas descritas a seguir em \texttt{platform/frontend/src/pages/}.

\subsection{Dashboard (Observabilidade)}
A p\'agina \textbf{Observabilidade} (\texttt{Dashboard.tsx}) \`e o ponto de entrada para inspe\c{c}\~ao do cat\'alogo e dos dados ao vivo.

\begin{itemize}
\item \textbf{Estado do sistema}: exibe o estado do SSE (\texttt{stream: open|connecting|closed}), o estado de carregamento das defini\c{c}\~oes (\texttt{/metrics/definitions}), contadores de eventos e timestamp da \''ultima amostra.
\item \textbf{Vistas e filtros}: organiza o cat\'alogo em vistas predefinidas (\textit{overview} e subconjuntos por p\'agina) e aplica filtros por camada, categoria, \textit{kind} e busca textual (nome/descri\c{c}\~ao).
\item \textbf{KPIs E2E}: no \textit{overview}, mostra KPIs de servi\c{c}o fim-a-fim (disponibilidade, lat\^encia, perda e jitter), calculados a partir das m\'etricas \texttt{e2e\_*} presentes no stream.
\item \textbf{Camadas}: apresenta, por camada, o total de defini\c{c}\~oes e quantas possuem dados observados no stream, permitindo identificar lacunas de instrumenta\c{c}\~ao.
\item \textbf{Piores caminhos}: monta um ranking \textit{Top 10} usando uma fun\c{c}\~ao de score baseada em perda e percentis (p95) de lat\^encia/jitter; o objetivo \`e priorizar diagn\'ostico de degrada\c{c}\~ao.
\end{itemize}

\subsection{Topologia}
A p\'agina \textbf{Topologia} (\texttt{Topology.tsx}) consome o mesmo SSE para exibir a topologia atual (hosts, switches e controlador). A visualiza\c{c}\~ao \`e propositalmente simplificada: agrupa n\'os por tipo e lista links com atributos (quando presentes), servindo como confirma\c{c}\~ao de invent\'ario e base de interpreta\c{c}\~ao para m\'etricas por n\'o e por caminho.

\subsection{Experimentos}
A p\'agina \textbf{Experimentos} (\texttt{Experiments.tsx}) estrutura experimenta\c{c}\~ao como registro cient\'ifico:

\begin{itemize}
\item \textbf{Templates (1 clique)}: biblioteca de perfis por categoria (ex.: L2, L3, L7, SDN, falhas, QoS), cada um definindo tr\'afego, dura\c{c}\~ao, camadas-alvo e um conjunto preferencial de m\'etricas. O usu\'ario pode \textit{usar} (preencher o formul\'ario) ou \textit{usar e executar} (criar e iniciar uma execu\c{c}\~ao imediatamente).
\item \textbf{Cria\c{c}\~ao manual}: permite especificar topologia, \texttt{src/dst}, gerador de tr\'afego (\texttt{ping}/\texttt{iperf}/\texttt{custom}), protocolo, taxa e dura\c{c}\~ao; e selecionar m\'etricas por camada a partir de \texttt{/metrics/definitions}.
\item \textbf{Execu\c{c}\~oes (runs)}: cada execu\c{c}\~ao \`e um \textit{snapshot} iniciado por \texttt{POST /experiments/\{id\}/run?mode=...}; ao final, os registros ficam dispon\'iveis em \texttt{GET /runs/\{run\_id\}/metrics}, com op\c{c}\~ao de exporta\c{c}\~ao JSON.
\item \textbf{Hist\'orico/resultados}: lista experimentos, execu\c{c}\~oes e apresenta uma tabela com os 50 registros mais recentes da execu\c{c}\~ao selecionada.
\end{itemize}

\subsection{Lab de Experimentos}
A p\'agina \textbf{Lab de Experimentos} (\texttt{Lab.tsx}) oferece uma constru\c{c}\~ao conceitual de \textit{workflow} por blocos (ex.: \textit{Ping}, \textit{TCP}, \textit{DNS}, \textit{Controlador SDN}, \textit{Fluxos OpenFlow}). Cada bloco declara: camada e m\'etricas-alvo; a tela filtra o stream para mostrar dados ao vivo correspondentes ao bloco selecionado. Trata-se de uma vista orientada a hip\'oteses e pipelines de medi\c{c}\~ao, n\~ao de execu\c{c}\~ao \textit{batch}.

\subsection{Monitor}
A p\'agina \textbf{Monitor} (\texttt{Monitor.tsx}) \`e a tela operacional de telemetria em tempo real via SSE:

\begin{itemize}
\item \textbf{Filtros por camada e n\'o}: permite selecionar subconjuntos para KPIs e para a tabela de eventos.
\item \textbf{KPIs ao vivo}: deriva indicadores simples a partir do primeiro valor dispon\'ivel entre m\'etricas alternativas (ex.: tr\'afego por \texttt{packet\_rate} ou \texttt{throughput\_mbps}; perda por \texttt{packet\_loss\_pct} etc.).
\item \textbf{Tabela de eventos}: mostra os 20 eventos mais recentes ap\'os filtros, com m\'etrica, n\'o, valor, camada e timestamp.
\end{itemize}

\subsection{Configura\c{c}\~oes (Settings)}
A p\'agina \textbf{Configura\c{c}\~oes} (\texttt{Settings.tsx}) \`e a principal interface para gest\~ao de perfis \texttt{platform\_config}:

\begin{itemize}
\item \textbf{Perfis salvos}: lista configs, permite ativar e remover (com prote\c{c}\~oes para perfil base).
\item \textbf{Edi\c{c}\~ao guiada}: monta uma configura\c{c}\~ao a partir de formul\'ario: modo (docker/local), controlador (opcional), n\'os (id/role/mgmt\_ip/container\_name), links (source/target e par\^ametros), camadas habilitadas e plano de m\'etricas (intervalo, labels e sele\c{c}\~ao por camada).
\item \textbf{Valida\c{c}\~ao}: imp\~oe restri\c{c}\~oes coerentes com a coleta (ex.: em docker, \texttt{container\_name}; para \texttt{ping}, hosts com \texttt{mgmt\_ip}; para \texttt{ovs-ofctl}, ao menos um switch).
\item \textbf{Preview JSON}: exibe o \texttt{platform\_config.json} derivado, servindo de ponte entre opera\c{c}\~ao assistida e reprodutibilidade baseada em arquivos.
\end{itemize}

\subsection{Hist\'orico}
A p\'agina \textbf{Hist\'orico} (\texttt{History.tsx}) encontra-se, no estado atual, como um \textit{placeholder} de UI (tabela de eventos simulados e bloco descritivo de relat\'orios). Do ponto de vista de produto, esta tela delimita uma extens\~ao natural: integra\c{c}\~ao com \texttt{/metrics/export} e com os arquivos \texttt{raw/*.jsonl} para recortes temporais, filtros e relat\'orios reprodut\'iveis.

\section{Discuss\~ao: extensibilidade e limita\c{c}\~oes}
A plataforma foi desenhada para evoluir sem quebrar contratos:

\begin{itemize}
\item \textbf{Extensibilidade por camadas}: o cat\'alogo estabelece um namespace e um vocabul\'ario (unidades, \textit{kind}, dimens\~oes) que permite adicionar novos coletores preservando compatibilidade com UI e exporta\c{c}\~ao.
\item \textbf{Reprodutibilidade}: a separa\c{c}\~ao entre telemetria ao vivo (SSE) e execu\c{c}\~oes registradas (runs + m\'etricas persistidas) endere\c{c}a simultaneamente opera\c{c}\~ao e pesquisa.
\item \textbf{Limita\c{c}\~oes atuais}: a tela de Hist\'orico \`e parcial; o modo real depende fortemente do ambiente (docker, permiss\~oes e ferramentas); e o score de \textit{piores caminhos} privilegia simplicidade operacional, podendo ser refinado com normaliza\c{c}\~ao por SLA/SLO e pesos por servi\c{c}o.
\end{itemize}

\section{S\'intese}
Este cap\'itulo apresentou a plataforma de observabilidade e experimenta\c{c}\~ao em SDN, destacando: (i) a arquitetura backend/frontend com SSE e REST; (ii) o modelo de dados e persist\^encia orientada a registros; (iii) o cat\'alogo completo de 267 m\'etricas e sua fun\c{c}\~ao como contrato; e (iv) o papel de cada p\'agina da interface (Dashboard, Topologia, Experimentos, Lab, Monitor, Configura\c{c}\~oes e Hist\'orico).

\end{document}